%=====================================================================
\chapter{Views, Triggers and Keeping Data Honest}\label{ch:views}
%=====================================================================
\locator{5}{Part V --- Getting the Design Right}

\begin{outcomes}
By the end of this chapter you will be able to:
\begin{tight}
  \item \textbf{Create} a view and \textbf{explain} what is stored when you
        do.
  \item \textbf{Show}, from a compiled program, that a view costs nothing at
        run time.
  \item \textbf{Use} a view to provide an external schema and to hide a
        schema change from the programs that read it.
  \item \textbf{Write} \sqlkw{BEFORE}, \sqlkw{AFTER} and \sqlkw{INSTEAD OF}
        triggers, and \textbf{say} when each is the right one.
  \item \textbf{Choose} between a view and a trigger-maintained table by what
        each guarantees.
\end{tight}
\end{outcomes}

\begin{prereq}
\chref{architecture}: the external level, which a view implements.
\chref{decomposition}: the summary table and its missing \sqlkw{UPDATE}
trigger. You need \texttt{code/m20\_views.py}.
\end{prereq}

\begin{keyterms}
view $\bullet$ virtual table $\bullet$ base table $\bullet$ view expansion
$\bullet$ flattening $\bullet$ co-routine $\bullet$ updatable view $\bullet$
trigger $\bullet$ \sqlkw{BEFORE} $\bullet$ \sqlkw{AFTER} $\bullet$
\sqlkw{INSTEAD OF} $\bullet$ \texttt{NEW} $\bullet$ \texttt{OLD} $\bullet$
\sqlkw{RAISE} $\bullet$ audit table $\bullet$ business rule $\bullet$
materialised view
\end{keyterms}

\section{A Query With a Name}

\begin{definitionbox}[View]
A \term{view} is a named query. \sqlkw{CREATE VIEW} $v$ \sqlkw{AS}
\sqlkw{SELECT} \dots\ stores the \emph{text} of the \sqlkw{SELECT} in the data
dictionary (\chref{whodoeswhat}) and nothing else --- no rows.

\smallskip
A query that names $v$ is rewritten, before it is compiled, with $v$'s
\sqlkw{SELECT} in its place. The view is a \term{virtual table}; the tables it
reads are \term{base tables}.

\smallskip
So a view is always \textbf{current}: it is the query, run now, against
whatever the base tables hold now. It cannot be stale, because there is
nothing stored to go stale.
\end{definitionbox}

\dsfig{%
\begin{tikzpicture}[
  q/.style={rounded corners=3pt, draw=#1, line width=1pt, fill=#1!7,
            text=#1!60!black, font=\scriptsize\ttfamily, align=left,
            text width=5.0cm, inner sep=4pt},
  dd/.style={rounded corners=3pt, draw=goldhl, line width=1pt, fill=goldhl!8,
             text=goldhl!60!black, font=\scriptsize\ttfamily, align=left,
             text width=4.4cm, inner sep=4pt},
  st/.style={rounded corners=3pt, draw=neutral, line width=0.9pt,
             fill=neutral!8, text=neutral!80!black, font=\scriptsize\bfseries,
             align=center, text width=2.6cm, minimum height=0.8cm},
  arr/.style={-{Stealth[length=2.2mm]}, primary!55, line width=0.9pt},
  tg/.style={font=\scriptsize\itshape, text=neutral}
]
  \node[q=secondary] (u) at (-4.4,1.4)
       {SELECT * FROM transcript\\ WHERE roll\_no = ?};
  \node[tg, anchor=south] at (-4.4,2.05) {what the program sends};
  \node[dd] (v) at (2.0,1.4)
       {transcript = SELECT e.roll\_no, c.course\_code, \dots\\
        FROM enrolment e JOIN course c ON \dots};
  \node[tg, anchor=south] at (2.0,2.2) {what the dictionary holds --- text};
  \node[q=primary] (r) at (-1.2,-0.4)
       {SELECT e.roll\_no, c.course\_code, \dots\\
        FROM enrolment e JOIN course c ON \dots\\
        WHERE e.roll\_no = ?};
  \node[tg, anchor=north, text width=5.4cm, align=center] at (-1.2,-1.25)
       {the query after substitution: the view is gone, and the filter sits on
        the base table};
  \node[st] (c) at (4.6,-0.4) {compiled\\program};
  \node[tg, anchor=north, text width=3.0cm, align=center] at (4.6,-0.95)
       {23 instructions, as if written by hand};
  \draw[arr] (u.south) -- ++(0,-0.35) -| ([xshift=-1.0cm]r.north);
  \draw[arr] (v.south) -- ++(0,-0.35) -| ([xshift=1.2cm]r.north);
  \draw[arr] (r.east) -- (c.west);
\end{tikzpicture}}%
{What a view is, in one diagram. The dictionary holds only the view's text;
SQLite substitutes it into the query that names it, then compiles. By the time
anything runs there is no view left to cost anything, or to be out of date.}%
{viewexpand}

\begin{examplebox}[The three jobs a view does in this course]
\textbf{1. An external schema.} \chref{architecture}'s \texttt{fees\_view} ---
roll number and name, nothing else --- is the external level, built:

\begin{lstlisting}[language=sql]
CREATE VIEW fees_view AS SELECT roll_no, name FROM student;
\end{lstlisting}

\textbf{2. Logical data independence.} When \chref{evolution}'s migration
split a column, a view with the old shape let old programs keep working:

\begin{lstlisting}[language=sql]
CREATE VIEW student_v1 AS
SELECT roll_no, name, programme, semester,
       addr_house || ', ' || addr_street || ', ' || addr_city AS address
  FROM student;
\end{lstlisting}

\textbf{3. A query written once.} The transcript is a join every report needs.
Write it once and every query that names \texttt{transcript} gets it right.

\smallskip
\textbf{What a view does \emph{not} do in SQLite} is restrict anyone to it.
There are no users (\chref{whodoeswhat}), so a person who can open the file can
read \textit{student} directly. The external level exists here; its
enforcement waits for \chref{project}.
\end{examplebox}

\section{Measured: What a View Costs, and What a Trigger Costs}

\begin{measured}[A view costs nothing; a trigger costs less than the same work by hand]
\textsf{\footnotesize \texttt{code/m20\_views.py}, seeded database (4{,}000
students), SQLite 3.51.2, Intel i5-6600, Windows 11, mechanical disk.}

\medskip
\textbf{1. A view against the query it wraps} --- one student's transcript

\begin{tabular}{@{}L{6.4cm}R{2.0cm}R{1.6cm}R{2.6cm}@{}}
& \textbf{ms} & \textbf{rows} & \textbf{instructions} \\
the query written out & 0.091 & 15 & 23 \\
the same through a view & 0.104 & 15 & 23 \\
\end{tabular}

\smallskip
\textbf{The same 23 instructions in the same order.} Seventeen of them differ
at all, and every one differs only in which \emph{cursor number} it uses: the
view's program numbers its tables one higher. That is the whole difference ---
the view was substituted and compiled away before it could cost anything.

\medskip
\textbf{2. Does a filter on the view reach inside it?}

\begin{tabular}{@{}L{7.0cm}L{7.0cm}@{}}
\texttt{transcript WHERE roll\_no = ?} &
  \texttt{SEARCH e USING INDEX \dots\ (roll\_no=?)} \\
grouped view, no index on course\_code &
  \texttt{CO-ROUTINE} + \texttt{SCAN enrolment} \\
grouped view, with that index &
  \texttt{CO-ROUTINE} + \texttt{SEARCH \dots\ (course\_code=?)} \\
\end{tabular}

\smallskip
A grouped view cannot be flattened, so it runs as a co-routine --- but the
filter is still pushed into it. Without an index the plan shows a
\texttt{SCAN}, which looks like no push-down and is not: there was simply
nothing to search with. With the index, only one course's rows were ever
grouped (0.277\,ms through the view, 0.268\,ms written out).

\medskip
\textbf{3. What a trigger costs} --- an audit trigger recording every grade
change; 200 changes, each done twice; median of 3 runs

\begin{tabular}{@{}L{6.2cm}R{2.6cm}R{2.4cm}R{2.6cm}@{}}
& \textbf{commit per row} & \textbf{(range)} & \textbf{one commit} \\
\sqlkw{UPDATE} only & 21.798\,ms & 19.4--22.8 & 0.1932\,ms \\
\sqlkw{UPDATE} + audit trigger & 22.222\,ms & 21.9--24.6 & 0.2234\,ms \\
\sqlkw{UPDATE} + the same \sqlkw{INSERT} by hand & 24.123\,ms & 21.6--24.1 &
  0.2810\,ms \\
\end{tabular}

\smallskip
\begin{tabular}{@{}L{6.2cm}R{3.6cm}R{4.0cm}@{}}
& \textbf{commit per row} & \textbf{one commit} \\
trigger adds per row & 0.423\,ms (+2\%) & 0.0302\,ms (+16\%) \\
the same work by hand adds & 2.324\,ms (+11\%) & 0.0879\,ms (+45\%) \\
\end{tabular}

\smallskip
\textbf{The trigger was cheaper than doing the same thing by hand} --- a third
of the overhead in bulk. It runs inside the \sqlkw{UPDATE} statement, in the
engine; the hand-written \sqlkw{INSERT} is a second statement the program must
send for every row. With a commit per row, the plain and trigger ranges
overlap: the disk sync dominates, and the trigger's share is inside its noise.

\smallskip
\textbf{A note on how this was established.} The first version of this
harness ran each case once, attributed the trigger's cost to ``the second
table'', and was wrong: a control written by hand cost \emph{more}, and single
runs of a 22\,ms commit on a mechanical disk vary by several milliseconds.
Medians, ranges and a control are what turned a guess into a result.

\medskip
\textbf{4. The view and \chref{decomposition}'s summary table, after one
enrolment moves from a CS course to an MA course}

\begin{tabular}{@{}L{3.0cm}R{2.4cm}R{2.4cm}R{2.8cm}@{}}
& \textbf{truth} & \textbf{view} & \textbf{summary table} \\
CS & 7{,}868 & 7{,}868 & \textbf{7{,}869} \\
MA & 7{,}174 & 7{,}174 & \textbf{7{,}173} \\
read cost & & 53.69\,ms & 0.076\,ms \\
\end{tabular}

\smallskip
The view is right, because it is the query run now. The table is wrong in both
rows, because it has insert and delete triggers and no update trigger --- the
gap \chref{decomposition} left. It is 700$\times$ faster to read and exactly as
correct as its triggers are complete.
\end{measured}

\begin{alertbox}[Choosing between a view and a maintained table]
\begin{tight}
  \item \textbf{Default to the view.} It costs nothing beyond the query it
        names, and it cannot be wrong. Measured.
  \item \textbf{Replace it with a maintained table only when the measured read
        cost forces you to} --- and then with triggers for \sqlkw{INSERT},
        \sqlkw{DELETE} \emph{and} \sqlkw{UPDATE}, plus a query that checks the
        table against the view, run on a schedule.
  \item \textbf{Do not fear triggers for their cost.} An audit trigger added
        16\% to an in-bulk update and was cheaper than the same work done by
        the program. What they cost is the attention needed to keep them
        complete.
  \item \textbf{Other systems offer a \term{materialised view}} --- a view
        whose result is stored and refreshed by the database. SQLite does not.
        The summary table with triggers is what that feature does, written out.
\end{tight}
\end{alertbox}

\section{Triggers}

\begin{definitionbox}[Trigger]
A \term{trigger} is a block of SQL the database runs automatically when a
named event happens to a table: \sqlkw{INSERT}, \sqlkw{UPDATE} (optionally
\sqlkw{OF} named columns) or \sqlkw{DELETE}.

\smallskip
It runs \sqlkw{BEFORE} or \sqlkw{AFTER} the change, once for each row
affected. Inside it, \texttt{NEW} is the row as it will be and \texttt{OLD} the
row as it was --- \texttt{NEW} has no meaning in a \sqlkw{DELETE} trigger, nor
\texttt{OLD} in an \sqlkw{INSERT} trigger.

\smallskip
\sqlkw{WHEN} limits it to rows meeting a condition. \texttt{RAISE(ABORT,
'message')} inside a trigger cancels the statement with an error.
\end{definitionbox}

\begin{worked}[A business rule no \sqlkw{CHECK} can express]
\textbf{The rule.} A student may not enrol in more than six sections in one
semester.

\medskip
\textbf{Step 1 --- why a \sqlkw{CHECK} cannot do it.} A \sqlkw{CHECK} sees one
row of one table. ``More than six'' is a fact about \emph{other} rows ---
the student's other enrolments --- so it is out of reach, exactly as
\chref{decomposition}'s two-table dependency was.

\medskip
\textbf{Step 2 --- \sqlkw{BEFORE INSERT}, so the bad row never arrives.}
\begin{lstlisting}[language=sql]
CREATE TRIGGER max_six BEFORE INSERT ON enrolment
WHEN (SELECT COUNT(*) FROM enrolment
       WHERE roll_no = NEW.roll_no AND semester = NEW.semester) >= 6
BEGIN
  SELECT RAISE(ABORT, 'a student may take at most six sections a term');
END;
\end{lstlisting}

\medskip
\textbf{Step 3 --- the gap.} This guards \sqlkw{INSERT}. An \sqlkw{UPDATE} that
moves an enrolment from Fall into a Spring that is already full is not
guarded. \textbf{Every rule implemented by trigger needs a trigger for every
event that can break it} --- the lesson of the summary table, again, and the
single most common defect in trigger code.

\medskip
\textbf{Step 4 --- why \sqlkw{BEFORE} and not \sqlkw{AFTER}.} Either would
work: \texttt{RAISE(ABORT)} in an \sqlkw{AFTER} trigger also undoes the
statement. \sqlkw{BEFORE} states the intent --- \emph{check, then allow} ---
and does no work on a row that will be refused.
\end{worked}

\begin{examplebox}[\sqlkw{INSTEAD OF}: making a view accept writes]
In SQLite \textbf{no} view can be written to directly --- not even
\texttt{fees\_view}, which reads a single table:
\texttt{cannot modify fees\_view because it is a view}. Other systems accept
writes to simple one-table views; SQLite accepts none, and an \sqlkw{INSTEAD
OF} trigger is the only way to say what a write to a view should do:

\begin{lstlisting}[language=sql]
CREATE TRIGGER fees_view_update INSTEAD OF UPDATE OF name ON fees_view
BEGIN
  UPDATE student SET name = NEW.name WHERE roll_no = OLD.roll_no;
END;

UPDATE fees_view SET name = 'Ayesha K. Khan'
 WHERE roll_no = '2021-IT-001';      -- now works
\end{lstlisting}

\smallskip
This completes the external level: a group of users can read \emph{and}
correct the part of the database they see, through the view, without ever
naming the base table.
\end{examplebox}

\begin{pitfall}
\begin{tight}
  \item \textbf{Believing a view stores rows.} It stores a \sqlkw{SELECT}. It
        cannot be stale.
  \item \textbf{Avoiding views for speed.} Measured: compiled away; same
        instructions as the query written out.
  \item \textbf{Reading \texttt{SCAN} in a grouped view's plan as ``no
        push-down''.} It may mean there was no index to search with.
  \item \textbf{A trigger for one event when the rule can be broken by three.}
        The summary table had no \sqlkw{UPDATE} trigger and was wrong after one
        \sqlkw{UPDATE}.
  \item \textbf{Using \texttt{NEW} in a \sqlkw{DELETE} trigger.} There is no
        new row.
  \item \textbf{Treating a view as security in SQLite.} No users, so no
        restriction.
  \item \textbf{Concluding from one timing run.} The first version of this
        chapter's harness did, and its explanation was wrong.
\end{tight}
\end{pitfall}

\begin{lab}[Prove a view is free, catch a summary table lying, and enforce a rule]
\begin{lstlisting}[language=bash]
# --- 1. The seeded database ----------------------------------------
python code/schema/seed.py university.db
sqlite3 university.db
\end{lstlisting}
\begin{lstlisting}[language=sql]
-- --- 2. A view, and the query it wraps --------------------------
CREATE VIEW transcript AS
SELECT e.roll_no, c.course_code, c.title, e.semester, e.grade
  FROM enrolment e JOIN course c ON c.course_code = e.course_code;
EXPLAIN QUERY PLAN SELECT * FROM transcript
 WHERE roll_no = (SELECT MIN(roll_no) FROM student);
-- The filter reached the base table. Where can you see that?

-- --- 3. What was stored -----------------------------------------
SELECT type, name, sql FROM sqlite_master WHERE name = 'transcript';
-- Text. No rows. That is the whole of a view.

-- --- 4. An external schema --------------------------------------
CREATE VIEW fees_view AS SELECT roll_no, name FROM student;
SELECT * FROM fees_view LIMIT 3;
SELECT * FROM student LIMIT 3;
-- The second still works. What would a server DBMS do instead?

-- --- 5. A rule no CHECK can express -----------------------------
CREATE TRIGGER max_six BEFORE INSERT ON enrolment
WHEN (SELECT COUNT(*) FROM enrolment WHERE roll_no = NEW.roll_no
                                       AND semester = NEW.semester) >= 6
BEGIN SELECT RAISE(ABORT, 'at most six sections a term'); END;
SELECT roll_no, COUNT(*) FROM enrolment WHERE semester = 'FALL2026'
 GROUP BY roll_no HAVING COUNT(*) = 6 LIMIT 1;
-- Try to enrol that student in one more FALL2026 section. Record
-- the message. Then find a way around the trigger with an UPDATE.

-- --- 6. An audit trail ------------------------------------------
CREATE TABLE grade_audit (roll_no TEXT, course_code TEXT,
  old_grade TEXT, new_grade TEXT,
  changed_at TEXT DEFAULT (datetime('now')));
CREATE TRIGGER audit_grade AFTER UPDATE OF grade ON enrolment
WHEN OLD.grade IS NOT NEW.grade BEGIN
  INSERT INTO grade_audit (roll_no, course_code, old_grade, new_grade)
  VALUES (OLD.roll_no, OLD.course_code, OLD.grade, NEW.grade);
END;
UPDATE enrolment SET grade = 'A'
 WHERE rowid = (SELECT MIN(rowid) FROM enrolment WHERE grade = 'B');
SELECT * FROM grade_audit;
\end{lstlisting}
\begin{lstlisting}[language=bash]
# --- 7. The measured version ----------------------------------------
python code/m20_views.py
\end{lstlisting}

\smallskip
\textbf{What to notice.} Step 5's second half is the chapter's warning, found
by hand: a rule enforced on \sqlkw{INSERT} can be broken by \sqlkw{UPDATE}.
Writing the \sqlkw{BEFORE UPDATE} trigger that closes the gap is part of the
project in \chref{project}.
\end{lab}

\begin{checkpoint}
\begin{tightnum}
  \item A view's compiled program had the same 23 instructions as the query it
        wraps, differing only in cursor numbers. Explain what SQLite did with
        the view to produce that, and why it means the view cannot be stale.
  \item The summary table read in 0.076\,ms and the view in 53.69\,ms, and
        after one \sqlkw{UPDATE} only the view was right. State the condition
        under which you would choose the table, and list everything you would
        have to build alongside it.
  \item Write the trigger that closes the gap in the six-section rule for
        \sqlkw{UPDATE}, and say which of \texttt{NEW} and \texttt{OLD} it must
        use, and why.
\end{tightnum}
\end{checkpoint}

\begin{chaptersummary}
\begin{tight}
  \item A \textbf{view} stores a \sqlkw{SELECT}, not rows. It is substituted
        into each query that names it, so it is always current.
  \item \textbf{Measured:} a view compiled to \textbf{the same 23
        instructions} as the query it wraps, differing only in cursor
        numbers. It costs nothing.
  \item \textbf{Measured:} a filter on a view reaches the base table, even
        through a grouped view run as a co-routine --- visible once there is an
        index to search with.
  \item Views provide the \textbf{external level}, \textbf{logical data
        independence}, and a query written once. In SQLite they do not restrict
        access.
  \item A \textbf{trigger} runs per row, \sqlkw{BEFORE} or \sqlkw{AFTER} an
        event, with \texttt{NEW} and \texttt{OLD}; \sqlkw{RAISE(ABORT)}
        refuses the change.
  \item \textbf{Measured:} an audit trigger added \textbf{16\%} in bulk and was
        cheaper than the same \sqlkw{INSERT} written by hand (45\%). With a
        commit per row its cost was inside the disk-sync noise.
  \item Triggers enforce rules a \sqlkw{CHECK} cannot reach --- but only for
        the events they are written for. Every event that can break the rule
        needs one.
  \item \textbf{Measured:} after one \sqlkw{UPDATE}, the view was right and
        the trigger-maintained table was wrong. Default to the view.
  \item \sqlkw{INSTEAD OF} triggers make a view writable.
\end{tight}
\end{chaptersummary}

\begin{reviewq}
  \item \textbf{(MCQ)} \sqlkw{CREATE VIEW} stores: (a)~a copy of the rows;
        (b)~the text of the query; (c)~an index; (d)~a trigger.
  \item \textbf{(MCQ)} In a \sqlkw{DELETE} trigger, which row reference is
        available? (a)~\texttt{NEW}; (b)~\texttt{OLD}; (c)~both;
        (d)~neither.
  \item \textbf{(Short)} Why can a view never be stale, while a summary table
        can?
  \item \textbf{(Short)} Give a rule that needs a trigger because no
        \sqlkw{CHECK} can express it, and say why.
  \item \textbf{(Short)} When is \sqlkw{INSTEAD OF} the only kind of trigger
        that will do?
  \item \textbf{(Applied)} Write the views that give the examinations office,
        the fees office and the student portal their external schemas from
        \chref{architecture}'s review question, an \sqlkw{INSTEAD OF} trigger
        letting the portal correct a student's phone number, and say what
        SQLite cannot enforce about any of it.
  \item \textbf{(Applied)} The audit trigger added 0.030\,ms per row in bulk.
        Estimate what auditing every grade change costs per semester on the
        seeded data, decide whether that matters, and list the three trigger
        events an audit of \textit{enrolment} would need to be complete.
\end{reviewq}
